\section{Adaptive estimation of the fidelity--0.5 boundary}

The goal is to estimate the set of randomized benchmarking parameters at which
the expected Boolean fidelity outcome is equal to one half.  In the present
experiments the two continuously varied parameters are the circuit depth \(d\)
and the minimum two-qubit gate ratio \(r\), with a fixed value of \(n\), the
number of qubits.  The target boundary is therefore
\[
    \Gamma_{0.5}
    =
    \{(d,r): p(d,r) = 0.5\},
\]
where \(p(d,r)\) denotes the probability that the backend returns a successful
fidelity outcome for a configuration with depth \(d\) and two-qubit ratio \(r\).
The methods in \texttt{viarregio4.py} and \texttt{viarregio5.py} estimate this
boundary under the prior structural assumption that fidelity is monotone:
increasing either \(d\) or \(r\) cannot improve the expected fidelity.

\subsection{Measurement model}

For each configuration \(x=(n,d,r)\), the backend is queried through
\texttt{fidelity\_estimation}.  Each measurement is converted to a Boolean
outcome
\[
    y_i(x) \in \{0,1\}.
\]
The empirical information at a configuration is summarized by the number of
successes \(s_x\) and failures \(f_x\).  A uniform Beta prior is used for the
unknown success probability \(p_x\),
\[
    p_x \sim \operatorname{Beta}(1,1).
\]
After the observations at \(x\), the posterior is
\[
    p_x \mid \{y_i(x)\}
    \sim
    \operatorname{Beta}(1+s_x, 1+f_x).
\]
The posterior mean
\[
    \widehat p_x
    =
    \mathbb{E}[p_x \mid \{y_i(x)\}]
    =
    \frac{s_x+1}{s_x+f_x+2}.
\]
is used when reporting or plotting the observed fidelity at a point.  The
surface fit itself uses the raw counts \(s_x\) and \(f_x\), as described below,
so that the uncertainty associated with low-shot points is handled directly by
the likelihood.

\subsection{HQC budget model}

The experimental budget is not simply the number of Boolean outcomes.  The
hardware quantum credit cost for running one circuit with shot count \(C\) is
modeled as
\[
    \operatorname{HQC}
    =
    5
    +
    \frac{N_1 + 10N_2 + 5N_m}{5000} C,
\]
where \(N_1\) is the number of one-qubit gates, \(N_2\) is the number of
two-qubit gates, \(N_m\) is the number of measurements, and \(C\) is the number
of times the same circuit is repeated.

The implementation uses this formula as the real budget when
\texttt{hqc\_budget} is set.  The older \texttt{measurement\_budget} remains as
a safety cap on the total number of Boolean outcomes, but the loop stops when
either the HQC budget or the measurement cap is exhausted.

Because the adaptive algorithm must choose configurations before the random
benchmarking circuit is explicitly generated, \texttt{viarregio4.py} uses
expected gate counts from the configuration.  For a configuration with \(n\)
qubits, depth \(d\), and two-qubit ratio \(r\), the random circuit has
\[
    nd
\]
gate slots before the inverse is appended.  A two-qubit gate occupies two of
these slots, so the expected number of two-qubit gates in the forward random
circuit is approximated by
\[
    N_{2,\mathrm{forward}}
    =
    \left\lfloor \frac{rnd}{2} \right\rfloor.
\]
The remaining occupied slots are one-qubit gates:
\[
    N_{1,\mathrm{forward}}
    =
    nd - 2N_{2,\mathrm{forward}}.
\]
Since the RMB circuit includes the inverse of the random circuit, the estimated
gate counts used for HQC accounting are
\[
    N_1 = 2N_{1,\mathrm{forward}} + N_{1,\mathrm{scrambler}},
    \qquad
    N_2 = 2N_{2,\mathrm{forward}},
    \qquad
    N_m = n.
\]
The scrambler contribution is included in expectation as
\[
    N_{1,\mathrm{scrambler}}
    =
    \operatorname{round}(2n p_{\mathrm{scramble}}),
\]
where \(p_{\mathrm{scramble}}\) is \texttt{scrambling\_probability}.  Identity
scrambler gates are not charged as physical one-qubit gates.

The fixed cost \(5\) is charged once for each selected configuration execution,
and the shot-dependent term scales with the number of requested shots for that
configuration.  This means that taking five shots at one selected configuration
is cheaper than treating those five shots as five separate circuit submissions,
because the fixed cost is paid once.

The HQC model is also used during candidate selection.  A candidate's usual
boundary-information score is multiplied by an efficiency factor
\[
    E_{\mathrm{HQC}}(x)
    =
    \left(
        \frac{\operatorname{HQC}(x,C_{\mathrm{plan}})}
             {\operatorname{HQC}_{\mathrm{base}}}
    \right)^{-\gamma},
\]
where \(C_{\mathrm{plan}}\) is the shot count the adaptive shot rule would
request for that candidate, \(\operatorname{HQC}_{\mathrm{base}}=5\), and
\(\gamma=\texttt{hqc\_cost\_power}\).  With the default
\(\gamma=1\), the candidate score is approximately information per HQC.  If
\(\gamma=0\), HQC cost is ignored during acquisition and is used only as a hard
affordability constraint.  Larger values of \(\gamma\) make the method more
aggressively prefer cheaper circuits.

This cost factor is applied to both global boundary candidates and
contour-following candidates.  When no fitted surface is available yet, random
fallback candidates and the initial design are ordered by estimated HQC cost so
that scarce budget is spent on affordable probes first.

\subsection{Monotone fidelity surface}

For each fixed \(n\), the observed points are fit with a monotone logistic
surface.  The physical coordinates \((d,r)\) are first scaled to
\((u,v) \in [0,1]^2\).  This scaling is not a physical assumption; it is a
numerical convenience.  Depth and two-qubit ratio have different units and very
different numerical ranges, so using them directly would make the optimizer
more sensitive to depth than to ratio simply because depth is the larger
number.  After scaling, a unit change in \(u\) means moving across the full
allowed depth range, while a unit change in \(v\) means moving across the full
allowed ratio range.

The model is
\[
    p(d,r)
    =
    \sigma\!\left(
        \alpha - \phi(u,v)^\top \beta
    \right),
\]
where \(\sigma(z) = 1/(1+\exp(-z))\), and
\[
    \phi(u,v)
    =
    \begin{bmatrix}
    u,& v,& u^2,& v^2,& uv,& u^3,& v^3,& u^2v,& uv^2
    \end{bmatrix}^{\top}.
\]
The function \(\sigma\) is the logistic function.  It converts any real number
into a value between zero and one, so it is a natural way to model a probability.
The quantity inside \(\sigma\),
\[
    z(d,r) = \alpha - \phi(u,v)^\top\beta,
\]
is called the logit of the probability.  Equivalently,
\[
    z = \operatorname{logit}(p)
    =
    \log\!\left(\frac{p}{1-p}\right).
\]
The logit is the log-odds of success.  It is useful because probabilities are
bounded, whereas logits live on the whole real line.  For example, \(p=0.5\)
corresponds to \(z=0\), \(p>0.5\) corresponds to \(z>0\), and \(p<0.5\)
corresponds to \(z<0\).  Therefore finding the fidelity--0.5 boundary is
equivalent to finding the zero level set of the fitted logit surface:
\[
    \Gamma_{0.5}
    =
    \{(d,r): z(d,r)=0\}.
\]

The intercept \(\alpha\) sets the overall easiness of the experiment.  If
\(\alpha\) is large and positive, the fitted probability near the origin is
high.  The vector \(\beta\) controls how the probability falls as depth and
two-qubit ratio increase.  The feature vector \(\phi(u,v)\) contains simple
polynomial terms:
\[
    u,\; v,\; u^2,\; v^2,\; uv,\; u^3,\; v^3,\; u^2v,\; uv^2.
\]
The linear terms \(u\) and \(v\) allow depth and two-qubit ratio to have separate
first-order effects.  The quadratic and cubic terms allow the transition to
curve rather than being a straight line.  The interaction terms \(uv\),
\(u^2v\), and \(uv^2\) allow the effect of depth to depend on the two-qubit gate
ratio, and vice versa.  This is important because a circuit can fail from many
moderately noisy operations, from fewer high-impact two-qubit operations, or
from a mixture of both.

This choice is deliberately more structured than a fully flexible interpolator,
but less restrictive than a straight-line fit.  It gives the boundary enough
freedom to bend while keeping the low-budget fit stable.

The constraint
\[
    \beta_j \geq 0
    \quad \text{for all } j
\]
forces the logit, and therefore the expected fidelity, to be non-increasing as
either \(d\) or \(r\) increases.  This is the main monotonicity exploitation in
\texttt{viarregio4.py}: points closer to the origin are assumed to be easier,
and points farther from the origin are assumed to be harder.

This constraint is easy to see from the fitted logit.  All features in
\(\phi(u,v)\) are non-negative on \([0,1]^2\).  If each coefficient
\(\beta_j\) is non-negative, then increasing \(u\) or \(v\) can only increase
\(\phi(u,v)^\top\beta\).  Because this term is subtracted from \(\alpha\), the
logit \(z(d,r)\) can only decrease.  Since the logistic function is increasing,
the probability \(p(d,r)=\sigma(z(d,r))\) can only decrease as well.

The parameters are obtained by minimizing the binomial negative log-likelihood,
\[
    \mathcal{L}(\alpha,\beta)
    =
    -\sum_x
    \left[
        s_x \log p_\theta(x)
        +
        f_x\log(1-p_\theta(x))
    \right]
    +
    \frac{\lambda}{2}\|\beta\|_2^2,
\]
subject to \(\beta_j\geq 0\).  Here \(s_x\) is the number of successful Boolean
outcomes at configuration \(x\), \(f_x\) is the number of failed outcomes, and
\(\lambda\) is the \texttt{monotone\_l2} regularization parameter.

Each term in this loss has a specific role.

\[
    p_\theta(x)
\]
is the probability predicted by the monotone surface at configuration \(x\).
The subscript \(\theta\) is shorthand for the fitted parameters
\(\theta=(\alpha,\beta)\).

\[
    s_x \log p_\theta(x)
    +
    f_x\log(1-p_\theta(x))
\]
is the binomial log-likelihood contribution from all repeated measurements at
that configuration.  Each success contributes \(\log p_\theta(x)\), and each
failure contributes \(\log(1-p_\theta(x))\).  A point with many repeated
measurements therefore contributes more information to the fit than a point
with only one or two repeated measurements.

This automatically accounts for unequal point errors.  If a point has only two
measurements, it contributes only two Bernoulli observations to the loss, so the
optimizer is free to move the fitted surface away from that noisy estimate if
the rest of the data disagree.  If a point has twenty measurements, it
contributes twenty Bernoulli observations and has a much stronger influence.
There is no separate hand-chosen weight: the number of successes and failures
is the weight.

This is also better than weighting only by a posterior mean.  Two points may
have the same posterior mean but very different uncertainty.  For example,
\(1\) success and \(1\) failure is much less informative than \(50\) successes
and \(50\) failures, even though both suggest a probability near \(0.5\).  The
binomial likelihood preserves this distinction.

The minus sign in front of the sum converts a likelihood maximization problem
into a loss minimization problem.  Maximizing likelihood and minimizing negative
log-likelihood are the same objective written in opposite directions.

The final term,
\[
    \frac{\lambda}{2}\|\beta\|_2^2,
\]
is an \(L^2\) regularization penalty.  It discourages unnecessarily large
coefficients, which would produce a very sharp transition surface.  This is
especially helpful at low measurement budgets, where a few random outcomes
could otherwise make the fitted boundary move too aggressively.  The intercept
\(\alpha\) is not penalized, because it controls the overall vertical placement
of the logit surface rather than the surface's shape.

The optimization is solved with bound-constrained L-BFGS-B.  The intercept
\(\alpha\) is unconstrained, while every component of \(\beta\) has the lower
bound \(0\).  The result is a smooth probability surface that respects the known
monotonicity of the backend and whose \(p=0.5\) contour can be extracted as the
estimated boundary.

\subsection{Initial measurements}

The first measurements are not placed on the four corners of the domain.  The
corners are usually low-information under the monotonicity assumption: the
lower-left corner is expected to pass, the upper-right corner is expected to
fail, and the other corners are often dominated by one large stressor.  Instead,
the initial design combines two ingredients.

First, a Latin hypercube design generates candidate points in the interior of
the allowed rectangle.  The margin from the edges is controlled by
\texttt{initial\_edge\_margin}.  Second, a small set of deliberate bracketing
points is added:
\[
    (d_{\mathrm{mid}}, r_{\mathrm{mid}}), \quad
    (d_{\mathrm{high}}, r_{\mathrm{mid}}), \quad
    (d_{\mathrm{mid}}, r_{\mathrm{high}}), \quad
    (d_{\mathrm{high}}, r_{\mathrm{low}}), \quad
    (d_{\mathrm{low}}, r_{\mathrm{high}}).
\]
These are intended to give the model early evidence on both sides of the
transition without spending measurements at the obvious corners.

The initial stage is budgeted.  The candidate initial design is first limited
by the measurement cap,
\[
    B_{\mathrm{init}}
    =
    \min\!\left(
        B - B_{\mathrm{reserve}},
        \lfloor \rho_{\mathrm{init}} B \rfloor
    \right),
\]
where \(B\) is the total measurement budget,
\(B_{\mathrm{reserve}}\) is \texttt{reserve\_adaptive\_measurements}, and
\(\rho_{\mathrm{init}}\) is \texttt{initial\_budget\_fraction}.  Each selected
initial configuration receives at most \texttt{initial\_shots\_per\_config}
measurements.  If \texttt{hqc\_budget} is set, each initial execution is then
checked against the remaining HQC budget before it is run.  Configurations that
cannot be afforded under the HQC formula are skipped, and the initial stage
stops once either budget is exhausted.  If there are too many candidate initial
configurations, they are thinned by farthest-point sampling, starting near the
center and then spreading out.

\subsection{Adaptive batches}

After the initial data are collected, the algorithm repeatedly proposes a small
batch of new configurations until the measurement budget is exhausted.  The
batch contains configurations, not necessarily one measurement each.  The
configured \texttt{batch\_size} is the maximum exploratory batch size; in the
current script the default is \(4\).  In \texttt{viarregio4.py} the actual batch
size can be adapted during the run.  Early in the experiment the full batch is
used, because the location of the boundary is poorly known and it is useful to
collect several diverse measurements before refitting.  Later, once the fitted
surface contains a contour and there are enough near-boundary anchors, the
runtime batch can shrink toward \texttt{min\_batch\_size}.  This lets the method
refit more frequently while it follows the line.

For each adaptive update, the data are grouped by \(n\).  For each group with at
least \texttt{min\_fit\_points} measured configurations, the monotone surface is
fit.  The batch proposal then mixes two kinds of candidates:
\[
    N_{\mathrm{contour}}
    =
    \operatorname{round}
    \left(
        \texttt{batch\_size}
        \times
        \texttt{contour\_follow\_fraction}
    \right),
\]
and
\[
    N_{\mathrm{global}}
    =
    \texttt{batch\_size} - N_{\mathrm{contour}}.
\]
With the initial defaults, half of the proposed configurations are intended to
follow the current boundary, while the other half are global boundary-seeking
proposals.  This split is useful because contour following is efficient once a
reliable boundary segment has been found, but global proposals protect against
the method locking onto an incorrect local segment too early.

When adaptive batching is enabled, the split also changes with the maturity of
the fit.  If no fitted \(p=0.5\) contour exists, the algorithm keeps the
configured exploratory behavior.  If a contour exists and the number of
near-boundary anchors is at least \texttt{contour\_min\_anchors}, the runtime
batch size is reduced to an intermediate value and the contour-following
fraction is increased.  If the run is past
\texttt{adaptive\_batch\_late\_budget\_fraction} of the total budget and the
anchor count is at least
\[
    \texttt{adaptive\_batch\_anchor\_multiplier}
    \times
    \texttt{contour\_min\_anchors},
\]
the algorithm switches to the late-stage behavior:
\[
    \texttt{batch\_size}_{\mathrm{runtime}}
    =
    \texttt{min\_batch\_size},
\]
with a larger contour-following fraction
\texttt{late\_contour\_follow\_fraction} and a smaller exploration weight
\texttt{late\_exploration\_weight}.  The purpose of this late-stage rule is to
move from ``find the boundary'' to ``trace the boundary''.

\subsubsection{Contour-following candidates}

Contour following begins from measured anchor points.  A configuration is an
anchor if the fitted surface predicts it to be near the boundary:
\[
    |p_\theta(d,r)-0.5|
    \leq
    \texttt{contour\_ready\_probability\_width}.
\]
Contour following is only enabled when at least
\texttt{contour\_min\_anchors} such anchors are present.

For each anchor, the algorithm estimates the local gradient
\[
    \nabla p_\theta(d,r)
\]
by finite differences.  This gradient is normal to the contour.  A tangent
direction is therefore
\[
    t = (\partial_r p_\theta, -\partial_d p_\theta).
\]
The method steps in both tangent directions by a fraction
\texttt{contour\_step\_fraction} of the parameter-range scale, clips the trial
point to the allowed rectangle, and then projects the trial point back toward
the fitted \(p=0.5\) contour by moving along the normal direction.  The
projection distance is limited by \texttt{contour\_projection\_fraction}, so the
candidate remains a local continuation of the known boundary rather than a
large jump.

Each contour-following candidate is scored by closeness to the fitted boundary
and by distance from existing measured points:
\[
    S_{\mathrm{contour}}(x)
    =
    \exp\!\left[
        -\left(
            \frac{|p_\theta(x)-0.5|}
                 {\texttt{boundary\_width}}
        \right)^2
    \right]
    \left(
        1
        +
        \texttt{exploration\_weight}
        \cdot
        \operatorname{dist}(x,\mathcal{D})
    \right).
\]
Configurations that have already reached \texttt{max\_shots\_per\_config} are
not proposed again.

\subsubsection{Global boundary candidates}

The global part of the batch optimizes a continuous acquisition function over
\((d,r)\).  Candidate points far from the fitted boundary are rejected:
\[
    |p_\theta(x)-0.5|
    >
    \texttt{max\_boundary\_probability\_distance}
    \quad \Rightarrow \quad
    S_{\mathrm{global}}(x)=0.
\]
Otherwise the score is
\[
    S_{\mathrm{global}}(x)
    =
    A(x)
    E(x)
    D(x)
    U(x),
\]
where
\[
    A(x)
    =
    \left[
    \exp\!\left(
        -\left(
            \frac{|p_\theta(x)-0.5|}
                 {\texttt{boundary\_width}}
        \right)^2
    \right)
    \right]^{\texttt{boundary\_focus\_power}}
\]
is the boundary-relevance term,
\[
    E(x)
    =
    1+\texttt{exploration\_weight}\cdot \operatorname{sparsity}(x)
\]
encourages sampling in parts of the boundary region with little previous data,
\[
    D(x)
    =
    \texttt{diversity\_floor}
    +
    (1-\texttt{diversity\_floor})\operatorname{diversity}(x)
\]
discourages proposing the same part of the parameter space repeatedly within
one batch, and
\[
    U(x)
    =
    1 - \frac{\text{current shots at }x}
             {\texttt{max\_shots\_per\_config}}
\]
discourages oversampling configurations that are already saturated.

This acquisition function is optimized with differential evolution.  The
continuous optimum is then rounded to an executable \texttt{RMBConfig}: depth is
rounded to an integer and the two-qubit ratio is rounded to two decimal places.

\subsection{Combining and executing a batch}

All contour-following and global candidates are pooled, sorted by score, and
deduplicated.  The top runtime batch size of unique configurations is selected.
If no fitted surface is available yet, because there are too few measured
points, the missing proposals are filled by random configurations from the
allowed parameter rectangle.

The batch size is therefore partly user-controlled and partly diagnostic-driven.
A batch size of one is the most greedy strategy: refit after every single
proposed configuration.  That can be measurement-efficient once the boundary is
already known locally, but it is computationally expensive and can be noisy at
very low budget, because one random Boolean outcome can move the next point too
strongly.  A larger early batch gives the model several pieces of evidence
before refitting, allows several circuit configurations to be prepared or run
together, and includes an explicit diversity term so the batch does not collapse
onto one rounded parameter value.  Very large batches would be less adaptive,
because many measurements would be chosen from a stale fitted surface.  The
adaptive rule uses the configured \texttt{batch\_size=4} for early exploration
and shrinks toward \texttt{min\_batch\_size=1} only when the observed data show
that the contour is mature enough to follow.

Once the configurations are chosen, the number of measurements assigned to each
configuration is adaptive.  If \(m_x\) measurements have already been taken at
\(x\), the configuration can receive at most
\[
    \texttt{max\_shots\_per\_config} - m_x
\]
additional measurements.  The requested number of new measurements is larger
when the current posterior mean is close to \(0.5\) and the posterior variance
is large:
\[
    R(x)
    =
    \exp\!\left[
        -\left(
            \frac{\widehat p_x-0.5}
                 {\texttt{shot\_boundary\_width}}
        \right)^2
    \right]
    \cdot
    \frac{\operatorname{Var}(p_x \mid \mathrm{data})}{1/12}.
\]
The value \(1/12\) is the variance of the initial
\(\operatorname{Beta}(1,1)\) prior.  The final shot count is clipped between
\texttt{min\_adaptive\_shots\_per\_config} and
\texttt{max\_adaptive\_shots\_per\_config}, by the per-configuration cap, by
the remaining measurement cap, and by the remaining HQC budget.  Under HQC
accounting, the algorithm computes the largest affordable shot count \(C\) for
the proposed configuration from
\[
    5
    +
    \frac{N_1 + 10N_2 + 5N_m}{5000} C
    \leq
    \operatorname{HQC}_{\mathrm{remaining}}.
\]
If even one shot cannot be afforded after paying the fixed cost, the
configuration is skipped.

\subsection{Current contour-first behaviour in \texttt{viarregio5.py}}

The current \texttt{viarregio5.py} strategy is more specialized than the
general adaptive batch strategy above.  It is designed around the belief that
the most efficient experiment is:
\[
    \text{find one reliable boundary anchor}
    \quad \longrightarrow \quad
    \text{trace the boundary}
    \quad \longrightarrow \quad
    \text{backfill uncertain traced points}.
\]
It therefore avoids repeatedly proposing global search batches once a plausible
contour has been found.

\subsubsection{Phase 1: protected initial crossing search}

The first task is to find one point on the \(p=0.5\) contour.  The method
starts at low two-qubit ratio because this is usually cheaper under the HQC
formula and often gives a useful entry point to the contour.  More generally,
it considers up to \texttt{ray\_ratio\_count} fixed-ratio rays between the
allowed minimum and maximum ratio, but it stops as soon as it finds a crossing.

For a chosen fixed ratio \(r\), the method searches in depth.  The current
implementation does not begin by probing the minimum and maximum depths.  Those
extreme probes are often low information: for example, at low two-qubit ratio a
zero- or near-zero-depth circuit is expected to pass and is not a credible
crossing point.  Instead, if
\texttt{initial\_crossing\_interior\_bracket} is enabled, the initial bracket is
placed in the interior of the depth range:
\[
    d_{\mathrm{center}}
    =
    d_{\min}
    +
    \texttt{initial\_crossing\_center\_fraction}
    (d_{\max}-d_{\min}),
\]
\[
    d_{\mathrm{low}}
    =
    d_{\mathrm{center}}
    -
    \texttt{initial\_crossing\_half\_width\_fraction}
    (d_{\max}-d_{\min}),
    \qquad
    d_{\mathrm{high}}
    =
    d_{\mathrm{center}}
    +
    \texttt{initial\_crossing\_half\_width\_fraction}
    (d_{\max}-d_{\min}).
\]
The bracket is clipped to the allowed depth range.  If the two endpoints do not
bracket a crossing, the bracket is expanded outward by a factor controlled by
\texttt{initial\_crossing\_expand\_factor}.  This gives the method a way to
recover if the interior guess is wrong, while still avoiding the naive
``start at the extremes'' behaviour.

Once a passing side and a failing side have been found, the method performs a
short bisection in depth.  Each bisection point receives
\texttt{ray\_bisection\_shots} repeated measurements.  If the measured
probability at a bisection point is near \(0.5\), the point can be topped up to
\texttt{crossing\_decision\_confirm\_shots} measurements before the bracket is
updated.  This is controlled by
\texttt{crossing\_decision\_probability\_width}.  The purpose is to avoid
moving the bracket too far because of one unlucky Boolean outcome.

At very low budgets, an ambiguous point can be dangerous: if it is treated as a
clear pass, the algorithm may overestimate the crossing depth and then trace
the wrong contour.  For this reason the current default
\texttt{crossing\_ambiguous\_as\_failure=True} treats still-ambiguous initial
crossing decisions conservatively.  This biases the initial anchor toward
shallower depth rather than allowing a single uncertain point to push the trace
too deep.

The initial search is also protected by a trace reserve.  A fraction of the HQC
budget is held back for contour following:
\[
    B_{\mathrm{reserve}}
    =
    \max\!\left(
        \texttt{trace\_reserve\_min\_hqc},
        \texttt{trace\_reserve\_budget\_fraction}
        \times
        B_{\mathrm{HQC}}
    \right),
\]
clipped so that it cannot exceed the total HQC budget.  The initial crossing
search and initial anchor refinement are not allowed to spend this reserve.
This is important for low-budget runs: the algorithm should not spend all of
the budget finding a single crossing and then have nothing left to trace the
contour.

After the first anchor is found, \texttt{initial\_anchor\_refine} optionally
performs a local fixed-ratio refinement around that anchor.  This uses
\texttt{initial\_anchor\_refine\_shots},
\texttt{initial\_anchor\_refine\_steps}, and
\texttt{initial\_anchor\_search\_fraction}.  The confirmed anchor can also be
topped up to at least \texttt{initial\_anchor\_min\_runs} measurements, subject
to the trace reserve and the per-configuration shot cap.

\subsubsection{Phase 2: projected contour tracing}

Once an anchor has been found, \texttt{viarregio5.py} traces the contour in the
positive ratio direction by default.  The current anchor is denoted
\[
    x_k = (d_k,r_k).
\]
The next ratio is chosen by a fixed step,
\[
    r_{k+1}
    =
    r_k
    +
    \texttt{trace\_ratio\_step\_fraction}
    (r_{\max}-r_{\min}),
\]
clipped to the allowed range.

If a monotone surface can be fit to the data, the algorithm estimates the local
gradient \(\nabla p_\theta(d_k,r_k)\).  The gradient is normal to the contour,
so a tangent direction is
\[
    t = (\partial_r p_\theta, -\partial_d p_\theta).
\]
The method takes a small step along this tangent and then projects the trial
point back toward the fitted \(p_\theta=0.5\) contour.  This produces a depth
guess for the next ratio.  If \texttt{trace\_enforce\_monotone\_depth} is
enabled, then moving to larger \(r\) is not allowed to increase the depth:
\[
    d_{k+1} \leq d_k.
\]
This encodes the expected monotonic shape of the boundary: if the two-qubit
gate ratio becomes more stressful, fewer gates should be needed to reach the
same failure probability.

The projected point is then measured with \texttt{trace\_shots} repeats.  If
the measured posterior mean is sufficiently close to \(0.5\),
\[
    |\widehat p_x - 0.5|
    \leq
    \texttt{trace\_accept\_probability\_width},
\]
it is accepted as the next anchor.  If it is not close enough, the method makes
up to \texttt{trace\_correction\_steps} local depth corrections.  A measured
point with \(\widehat p_x>0.5\) is too easy, so depth is increased; a point with
\(\widehat p_x<0.5\) is too hard, so depth is decreased.  Each correction step
shrinks the depth step size, so the correction behaves like a short local
search rather than a large vertical sweep.

The current implementation also has a rejection rule.  If the best corrected
point is still farther than
\texttt{trace\_reject\_probability\_width} from \(0.5\), it is not accepted as
a new anchor and tracing stops.  In addition, if a candidate lies on the depth
boundary and is not already within
\texttt{trace\_accept\_probability\_width}, it is rejected.  This prevents the
algorithm from marching along a depth bound and producing wasteful vertical
bands of points that are not actually on the fidelity contour.

\subsubsection{Phase 3: contour-anchor backfill}

After tracing reaches the parameter bounds or stops because the next point is
not credible, the algorithm may still have HQC budget left.  The current
default does not use this remaining budget to start a new broad search.
Instead, if \texttt{refine\_after\_trace} is enabled, it first revisits the
anchors already traced along the contour.

Each traced anchor can be repeated until it reaches
\[
    \min(
        \texttt{trace\_anchor\_min\_shots},
        \texttt{max\_shots\_per\_config}
    )
\]
measurements.  The next anchor to repeat is chosen by posterior variance, so
the least certain traced points are backfilled first.  This increases
confidence along the discovered contour without introducing new fixed-depth or
fixed-ratio exploration bands.

If traced anchors have all reached their target repeat count and budget remains,
the method can then call the more general
\texttt{refine\_uncertain\_boundary\_points} rule.  This rule considers measured
configurations that are not yet saturated and scores them by two factors:
posterior variance and closeness of the fitted monotone surface to \(0.5\).  It
therefore spends leftover measurements on uncertain measured points near the
current fitted boundary.

\subsubsection{Key contour-search parameters}

The following parameters have the largest effect on the current contour search.

\begin{itemize}
    \item \texttt{ray\_ratio\_count}: maximum number of fixed-ratio rays tried
    while looking for the first contour anchor.  Larger values give more chances
    to find a good entry point, but can spend more budget before tracing begins.

    \item \texttt{ray\_probe\_shots}: number of repeats used for the initial
    low/high depth probes on a fixed-ratio ray.  Larger values make the bracket
    more reliable but more expensive.

    \item \texttt{ray\_bisection\_steps}: number of depth bisection steps used
    after a bracket is found.  More steps localize the initial crossing more
    precisely, but can leave less budget for contour following.

    \item \texttt{ray\_bisection\_shots}: repeats per initial bisection point.
    Increasing this reduces noise in the initial crossing estimate.

    \item \texttt{initial\_crossing\_interior\_bracket}: if enabled, the first
    fixed-ratio search starts from an interior depth bracket rather than the
    minimum and maximum depths.

    \item \texttt{initial\_crossing\_center\_fraction}: location of the interior
    bracket center as a fraction of the depth range.  Smaller values start the
    search at shallower depths; larger values start deeper.

    \item \texttt{initial\_crossing\_half\_width\_fraction}: half-width of the
    initial depth bracket as a fraction of the depth range.  Larger values make
    the initial bracket broader but may reintroduce low-information probes.

    \item \texttt{initial\_crossing\_expand\_factor}: controls how aggressively
    the interior bracket expands if it does not initially contain a crossing.

    \item \texttt{crossing\_decision\_confirm\_shots}: target number of repeats
    for an initial bisection point whose measured probability is near \(0.5\).
    This is used to avoid updating the bracket from a very noisy near-boundary
    measurement.

    \item \texttt{crossing\_decision\_probability\_width}: defines what
    ``near \(0.5\)'' means for the initial crossing decision.  Larger values
    confirm more points; smaller values make the algorithm cheaper but more
    reactive to noise.

    \item \texttt{crossing\_ambiguous\_as\_failure}: when enabled, ambiguous
    initial bisection points are treated conservatively after confirmation.
    This helps avoid overestimating the initial crossing depth.

    \item \texttt{initial\_anchor\_refine}: enables a local refinement pass at
    the first found anchor before contour tracing begins.

    \item \texttt{initial\_anchor\_refine\_shots}: repeats per point in the
    local initial-anchor refinement.

    \item \texttt{initial\_anchor\_refine\_steps}: number of local refinement
    steps around the first anchor.

    \item \texttt{initial\_anchor\_search\_fraction}: depth range used for local
    initial-anchor refinement, as a fraction of the full depth range.

    \item \texttt{initial\_anchor\_min\_runs}: minimum repeat count requested
    for the confirmed first anchor, subject to budget and shot caps.

    \item \texttt{trace\_reserve\_budget\_fraction}: fraction of the HQC budget
    reserved for tracing and not available to the initial crossing search.

    \item \texttt{trace\_reserve\_min\_hqc}: absolute minimum HQC reserve for
    tracing.  This is especially important for small total budgets.

    \item \texttt{trace\_ratio\_step\_fraction}: ratio step size used while
    following the contour.  Larger values traverse the contour faster but are
    more likely to miss curvature; smaller values are safer but require more
    circuit executions.

    \item \texttt{trace\_shots}: repeats used for each proposed traced point.
    Larger values give more reliable accept/reject decisions.

    \item \texttt{trace\_correction\_steps}: maximum number of local depth
    corrections after a projected trace point is measured off the boundary.
    Larger values can recover from projection error, but too many corrections
    can create vertical-looking local search patterns.

    \item \texttt{trace\_depth\_search\_fraction}: initial correction step size
    in depth, as a fraction of the full depth range.

    \item \texttt{trace\_accept\_probability\_width}: probability distance from
    \(0.5\) required to immediately accept a measured trace point as an anchor.

    \item \texttt{trace\_reject\_probability\_width}: maximum allowed distance
    from \(0.5\) after correction.  If the best candidate is still farther away,
    tracing stops instead of accepting a bad anchor.

    \item \texttt{trace\_enforce\_monotone\_depth}: prevents depth from
    increasing when the trace moves to larger two-qubit ratio.  This uses the
    expected monotone shape of the boundary directly.

    \item \texttt{trace\_anchor\_min\_shots}: target repeat count for backfilling
    already traced anchors after the contour-following phase.

    \item \texttt{refine\_after\_trace}: enables the post-trace backfill and
    remaining-boundary refinement stages.

    \item \texttt{refinement\_shots}: repeats added in each post-trace
    refinement execution.

    \item \texttt{refinement\_boundary\_width}: controls how close the fitted
    monotone surface must be to \(0.5\) for a measured point to be considered
    useful for leftover-budget refinement.

    \item \texttt{max\_shots\_per\_config}: maximum repeats allowed at any one
    configuration.  This bounds how much budget can be spent backfilling a
    single anchor.

    \item \texttt{hqc\_cost\_power}: strength of HQC cost-awareness in candidate
    scoring.  Larger values prefer cheaper circuits more strongly.
\end{itemize}

\subsection{Cost-aware stitched batching in \texttt{viarregio7.py}}

\texttt{viarregio7.py} keeps the contour-first logic of
\texttt{viarregio5.py}, but changes the execution model used for HQC
accounting.  The local experiment is still run through the \texttt{SympleQ}
backend.  The batching layer is only a cost model for the corresponding
Quantinuum-style execution: it asks how much the same set of local RMB circuits
would cost if several result circuits were stitched into one larger submitted
job.

The reason batching is useful is the flat term in the HQC formula.  If \(J\)
circuits are submitted independently, the flat contribution is
\[
    5J.
\]
If those same result circuits can be stitched into one homogeneous batch, the
flat contribution is instead
\[
    5.
\]
The gate and measurement terms do not disappear.  They are accumulated over all
result circuits in the stitched job, and a reset overhead is added between
subcircuits.  For a batch \(B\), the estimated stitched cost is
\[
    \operatorname{HQC}_{\mathrm{batch}}(B)
    =
    5
    +
    \frac{
        \sum_{x\in B} C_x W(x)
        +
        W_{\mathrm{reset}} n \left(\sum_{x\in B} C_x - 1\right)
    }{5000},
\]
where \(C_x\) is the requested shot count for configuration \(x\), \(W(x)\) is
the weighted one-qubit, two-qubit, and measurement size of one RMB result
circuit, and \(n\) is the common qubit count for the batch.  The reset term is
charged between stitched subcircuits.  The native unbatched comparison used in
diagnostics is
\[
    \operatorname{HQC}_{\mathrm{native}}
    =
    \sum_x C_x
    \left[
        5 + \frac{W(x)}{5000}
    \right].
\]
The reported batching saving is the difference between these two estimates.

Batches are kept homogeneous in \(n\).  This is important because a stitched
job with mixed qubit counts would not correspond to a simple single-register
execution model.  The batch packer therefore groups only configurations with
the same number of qubits, and it refuses to add a request if doing so would
exceed either the remaining global HQC budget or the configured
\texttt{max\_cost\_per\_batch}.  The parameter
\texttt{batch\_max\_configs} also limits the number of distinct configurations
in one stitched job.

The batcher tries to make each job large enough to benefit from the flat-cost
saving.  Its target is controlled by \texttt{batch\_target\_fill\_fraction};
with the current defaults, a job is packed toward a large fraction of
\texttt{max\_cost\_per\_batch} rather than being submitted as soon as a single
configuration is available.  If \texttt{batch\_fill\_repeats} is enabled, the
packer can also add extra repeats of configurations already in the batch until
the target fill is approached, subject to the per-configuration shot cap.  This
exploits the same flat-cost effect in another way: repeating a stitched set of
circuits is often cheaper than submitting many small jobs with the same total
number of Boolean outcomes.

The adaptive stages are still logically the same as before.  There is an
initial crossing search, then contour following, then refinement or backfill.
What changes is the unit of execution.  Instead of every proposed point being
costed as a separate job, the requests from a stage are packed into one or more
stitched jobs.  The data recorded after execution are still ordinary counts
\((s_x,f_x)\) at individual RMB configurations, so the statistical model does
not need to know whether those measurements came from a singleton job or a
stitched batch.  Batch identifiers are stored for diagnostics and plotting, so
one can see which measured points were chosen together.

\subsection{Three-dimensional extension in \texttt{viarregio8.py}}

\texttt{viarregio8.py} extends the problem from a two-dimensional contour at a
fixed qubit count to a three-dimensional fidelity volume over
\[
    x = (d,r,n),
\]
where \(d\) is the legacy depth coordinate, \(r\) is the minimum two-qubit gate
ratio, and \(n\) is the number of qubits.  The target is no longer just a curve.
It is the surface
\[
    \Gamma_{0.5}
    =
    \{(d,r,n): p(d,r,n)=0.5\}.
\]
The main scalar output is the normalized volume fraction of the configured
domain where the fitted probability is at least one half.

\subsubsection{Coupled monotone volume model}

The 3D model uses one coupled monotone logistic fit rather than fitting each
qubit slice independently.  Physical coordinates are scaled to
\[
    (u,v,w) \in [0,1]^3,
\]
corresponding to depth, ratio, and qubit count.  The fitted probability is
\[
    p(d,r,n)
    =
    \sigma\!\left(
        \alpha - \phi(u,v,w)^\top\beta
    \right),
\]
where the feature vector contains linear, quadratic, cubic, and interaction
terms,
\[
\begin{aligned}
    \phi(u,v,w) =
    (&u,\;v,\;w,\;u^2,\;v^2,\;w^2,\;uv,\;uw,\;vw,\\
     &u^3,\;v^3,\;w^3,\;uvw).
\end{aligned}
\]
As in the 2D case, the coefficients are constrained by
\[
    \beta_j \geq 0.
\]
This makes the fitted fidelity non-increasing as any of depth, two-qubit ratio,
or qubit count increases.  The coupled fit lets measurements at, for example,
\(n=30\) inform the expected surface at \(n=25\) and \(n=35\), while still
allowing curvature and interactions between the three axes.

The fit uses the same binomial likelihood style as the 2D monotone surface.
For each measured configuration, the likelihood contribution is based on its
observed successes and failures, not only on the posterior mean.  Thus a point
with many repeats has more influence than a point with one noisy Boolean
outcome.

\subsubsection{Qubit-dependent depth domain}

The 3D benchmark uses a qubit-dependent maximum depth.  A single global maximum
depth is inappropriate: a 50-qubit circuit may reach the \(p=0.5\) transition
at low depth, while a 10-qubit circuit can require a much larger depth.  The
current default depth domain is
\[
    d_{\max}(n)
    =
    \min\!\left(
        150,\;
        50\left(\frac{30}{n}\right)
    \right),
\]
with the minimum depth still set by \texttt{depth\_bounds}.  This gives
approximately
\[
    d_{\max}(10)\approx 150,\qquad
    d_{\max}(30)\approx 50,\qquad
    d_{\max}(50)\approx 30.
\]
The cap is controlled by \texttt{qubit\_depth\_cap\_enabled},
\texttt{qubit\_depth\_cap\_plateau},
\texttt{qubit\_depth\_cap\_power}, and
\texttt{qubit\_depth\_cap\_margin\_fraction}.  This qubit-dependent maximum is
used during seeding, acquisition, contour extraction, plotting, and normalized
volume computation.

\subsubsection{Height-field contour representation}

The 3D contour is represented as a depth height-field over the
\((r,n)\)-plane.  For each grid point \((r_j,n_k)\), the fitted model is queried
along the allowed depth interval
\[
    [d_{\min}, d_{\max}(n_k)].
\]
If the fitted probability crosses one half, the crossing depth
\[
    d_{50}(r_j,n_k)
\]
is found by bisection.  If the whole interval is above one half, the height is
set to \(d_{\max}(n_k)\), meaning the configured depth range is still inside the
passing region.  If the whole interval is below one half, the height is set to
\(d_{\min}\), meaning the passing volume has already disappeared at that
\((r,n)\) location.

The normalized volume is then estimated by averaging the normalized heights,
\[
    V_{\mathrm{norm}}
    =
    \frac{1}{N_rN_n}
    \sum_{j,k}
    \frac{
        d_{50}(r_j,n_k)-d_{\min}
    }{
        d_{\max}(n_k)-d_{\min}
    }.
\]
This number lies in \([0,1]\).  It is a normalized fraction of the configured
3D box, not a raw physical volume.  The saved volume grid contains the ratio
grid, qubit grid, \(d_{50}\) values, per-qubit \(d_{\max}\) values, crossing
statuses, and the normalized volume.

\subsubsection{3D acquisition}

The initial stage seeds a small number of qubit slices, currently
\texttt{seed\_n\_qubits\_values=(10,30,50)} by default.  Within each seed slice
the method uses the same lesson as the 2D contour-first algorithm.  It first
spends a batched low-ratio horizontal search in depth, using configured
\texttt{initial\_low\_ratio\_depth\_fractions}.  Only after those measurements
have been recorded does it estimate the low-ratio anchor depth and fan out in
ratio.  This makes the first contour seed depend on the measured crossing,
rather than on a fixed prior guess.  The initial design is therefore not a
uniform 3D grid.  It is a small number of contour-seeking probes in
representative qubit slices.

Once enough data exist, the coupled volume model is fit and used to propose new
locations.  Candidate \((r,n)\) pairs are drawn from a ratio-qubit grid.  At
each candidate pair, the fitted height-field gives a predicted depth
\(d_{50}(r,n)\).  The algorithm then asks for a same-ratio depth bracket around
that predicted crossing, using the qubit-dependent depth range for that slice.
The intended pattern is the 3D analogue of good 2D contour following: for most
sampled \((r,n)\) locations, the data should include points just above and just
below the transition.  This is important statistically: for a fixed \((r,n)\),
a few depths on both sides of the transition give the monotone model much
stronger evidence than many repeats at a point far from the boundary.

Before proposing broad new volume candidates, the current v8 method also runs a
bracket-completion pass.  This pass scans the measured data for fixed
\((r,n)\)-like groups that are still one-sided: all measured depths pass, or all
measured depths fail.  If such a group is close to the fitted contour, the
method asks for the missing side.  For an all-passing group it proposes a
deeper point; for an all-failing group it proposes a shallower point.  The
purpose is to make the 3D acquisition behave like the best 2D runs: after the
initial crossing is found, most followed ratios should have evidence on both
sides of the \(p=0.5\) transition.

Candidate scoring combines four pressures.  The first is contour uncertainty,
which is large near \(p=0.5\).  The second is sparsity in scaled
\((d,r,n)\)-space, which discourages repeatedly measuring the same part of the
volume.  The third is within-batch diversity, so one acquisition pass covers
several regions rather than collapsing onto one rounded configuration.  The
fourth is HQC efficiency, controlled by \texttt{volume\_cost\_power}.  Large,
high-qubit, high-ratio circuits are therefore not forbidden, but they must earn
their cost by reducing uncertainty in a part of the volume that matters.

The 3D acquisition requests are executed by the same stitched batch machinery
as v7.  Batches remain homogeneous in \(n\), because stitching mixed-qubit
circuits would not match the intended hardware-cost model.  The batch cap
\texttt{max\_cost\_per\_batch} is still enforced for every stitched job, and the
global HQC budget is checked before each batch is run.  Diagnostics report the
stitched HQC spent, the native unbatched estimate, the estimated saving, the
number of batched jobs, and the largest batched job size.

\subsection{Boundary extraction and uncertainty}

After the measurements are complete, the monotone surface is evaluated on a
regular plotting grid and the \(p_\theta(d,r)=0.5\) contour is extracted.  The
grid is used only for visualization and contour extraction; measurements are
not restricted to the grid during the adaptive stage.

Uncertainty in the boundary is estimated by posterior bootstrap.  For every
measured configuration \(x\), a sampled fidelity probability is drawn from
\[
    p_x^{(b)}
    \sim
    \operatorname{Beta}(1+s_x,1+f_x).
\]
For each bootstrap draw \(b\), a new monotone surface is fit to the sampled
probabilities, and its \(p=0.5\) contour is extracted.  The collection of
bootstrap contours approximates the posterior uncertainty over the boundary.
This bootstrap step deliberately uses the full Beta posterior at each measured
configuration, while the main fitted surface uses the binomial likelihood of
the observed counts.  Thus the central fit is count-likelihood based, and the
uncertainty visualization propagates the remaining posterior uncertainty in
those counts.
The plot shows this uncertainty as a contour-occupancy heatmap by default,
rather than as many individual contour lines, because broad uncertainty can
otherwise produce a visually cluttered set of vertical or crossing line
segments.

\subsection{Complete pipeline}

The full \texttt{viarregio4.py} pipeline is as follows.

\begin{enumerate}
    \item Define the total measurement budget and parameter bounds for
    \((n,d,r)\).

    \item Construct the \texttt{SympleqBackend} with the chosen one-qubit and
    two-qubit Pauli noise models.

    \item Generate an initial, budget-limited design.  The design uses interior
    Latin hypercube points plus non-corner bracketing points, then thins the
    candidates if the initial budget is too small.

    \item Spend the initial measurements.  Each Boolean outcome updates the
    \(\operatorname{Beta}(1+s_x,1+f_x)\) posterior for its configuration.

    \item Fit a monotone logistic fidelity surface separately for each fixed
    \(n\), provided at least \texttt{min\_fit\_points} configurations have been
    measured.

    \item Propose the next batch.  The runtime batch size and
    contour-following fraction are first chosen from the current maturity
    diagnostics: whether a contour exists, how many near-boundary anchors are
    available, and how much of the total budget has been spent.  A fraction of
    the batch then attempts to follow the currently fitted contour from
    near-boundary anchors.  The remainder is chosen by global optimization of
    the boundary-focused acquisition score.

    \item Deduplicate the candidate configurations, sort by score, and keep at
    most \texttt{batch\_size} configurations.

    \item For each selected configuration, choose an adaptive number of
    measurements.  More measurements are spent near the boundary and where the
    Beta posterior variance remains high.

    \item Record the outcomes, decrement the remaining global budget, and refit
    the monotone surface before selecting the next batch.

    \item When the budget is exhausted, extract the fitted \(p=0.5\) contour and
    estimate its uncertainty by posterior bootstrap.
\end{enumerate}

The current \texttt{viarregio5.py} pipeline is:

\begin{enumerate}
    \item Define the HQC budget, measurement cap, parameter bounds, and fixed
    \(n\)-qubit value.

    \item Reserve part of the HQC budget for contour following using
    \texttt{trace\_reserve\_budget\_fraction} and
    \texttt{trace\_reserve\_min\_hqc}.

    \item Try fixed-ratio depth searches, starting at low ratio and stopping
    when the first credible \(p=0.5\) crossing is found.

    \item For each fixed-ratio search, start from an interior depth bracket,
    expand only if necessary, and then run a short depth bisection.  Near-boundary
    bisection decisions can be confirmed with extra repeats before the bracket
    is updated.

    \item Optionally refine and top up the first anchor, without spending the
    reserved trace budget.

    \item Trace the contour in ratio.  At each step, use the fitted monotone
    surface to propose a projected next point, measure it, and make a small
    number of local depth corrections if the measured point is off the boundary.

    \item Reject a proposed trace point if it is too far from \(0.5\), or if it
    lies on a depth bound without being sufficiently close to \(0.5\).  This
    prevents the algorithm from turning a failed trace step into a vertical
    edge search.

    \item When tracing stops, backfill the traced anchors with additional
    repeats, prioritizing anchors with high posterior variance.

    \item If budget remains after anchor backfill, refine uncertain measured
    points that the fitted surface predicts to be near the boundary.

    \item Fit the final monotone surface, extract its \(p=0.5\) contour, and
    optionally estimate uncertainty by posterior bootstrap.
\end{enumerate}

The \texttt{viarregio7.py} batched contour-first pipeline is:

\begin{enumerate}
    \item Define the same fixed-\(n\) contour problem as v5, together with an
    overall HQC budget and \texttt{max\_cost\_per\_batch}.

    \item Build only the local \texttt{SympleQ} backend.  Quantinuum-style HQC
    formulas are used for costing and diagnostics, not for execution.

    \item Generate the initial crossing-search requests and pack them into
    homogeneous stitched batches where possible.  Each stitched batch pays the
    flat HQC base cost once.

    \item Run the local SympleQ circuits in the batch and record ordinary
    per-configuration Boolean counts.

    \item Trace the contour as in v5, but execute the proposed trace brackets
    and corrections through the same batch packer.

    \item Use leftover budget for anchor backfill or uncertain-boundary
    refinement, again packing requests toward the per-batch cost cap.

    \item Fit the final 2D monotone surface, extract the contour, and report
    both stitched HQC spent and the native unbatched HQC estimate.
\end{enumerate}

The \texttt{viarregio8.py} 3D volume pipeline is:

\begin{enumerate}
    \item Define the 3D domain in \((d,r,n)\), including the allowed
    \texttt{n\_qubits\_values} and the qubit-dependent depth maximum
    \(d_{\max}(n)\).

    \item Seed representative qubit slices, usually \(n=10,30,50\), with
    low-ratio depth searches and contour-bracket probes.

    \item Pack these seed requests into homogeneous stitched batches and run
    them locally through SympleQ.

    \item Fit the coupled monotone volume model once the measured data contain
    enough distinct configurations.

    \item Extract the current predicted height-field \(d_{50}(r,n)\), select
    new ratio-qubit locations near the predicted transition, and request small
    depth brackets around those predicted crossing depths.

    \item Score and pack the new requests using uncertainty, sparsity,
    diversity, and HQC efficiency, while respecting both the global budget and
    \texttt{max\_cost\_per\_batch}.

    \item Repeat the fit--acquire--batch loop until the acquisition pass limit
    or the HQC budget is exhausted.

    \item Fit the final coupled 3D model, save the \(d_{50}(r,n)\) grid, compute
    the normalized fidelity volume, and plot the 3D height-field with measured
    points overlaid.
\end{enumerate}

The essential idea is that monotonicity converts the problem from learning an
arbitrary surface into learning a structured transition boundary.  In two
dimensions this boundary is a curve; in three dimensions it is a height-field
surface whose integral gives the normalized passing volume.  The batching layer
is orthogonal to the statistical model: it changes how candidate circuits are
costed and grouped, while the fitted model still sees the resulting successes
and failures at individual RMB configurations.
